% Cabeçalho ABNT para o PDF (XeLaTeX + pandoc)
\usepackage{indentfirst}
\usepackage{titlesec}
\usepackage{fancyhdr}

\setlength{\parindent}{1.25cm}
\setlength{\parskip}{0pt}

\titleformat{\section}{\normalfont\bfseries}{\thesection}{0.75em}{}
\titlespacing*{\section}{0pt}{1.5\baselineskip}{1\baselineskip}
\titleformat{\subsection}{\normalfont\bfseries}{\thesubsection}{0.75em}{}
\titlespacing*{\subsection}{0pt}{1.2\baselineskip}{0.8\baselineskip}

\pagestyle{fancy}
\fancyhf{}
\fancyhead[R]{\thepage}
\renewcommand{\headrulewidth}{0pt}

% Citação longa: 4 cm de recuo, 10 pt, espaço simples.
\renewenvironment{quote}
  {\begin{list}{}{\leftmargin=4cm \rightmargin=0cm}\item\relax\small\setstretch{1.0}}
  {\end{list}}

% Resumo/abstract: sem recuo, espaço simples.
\newenvironment{resumo}
  {\par\setlength{\parindent}{0pt}\setstretch{1.0}}
  {\par}

% Referências: sem recuo, espaço simples e uma linha entre entradas.
\newenvironment{refs}
  {\par\setlength{\parindent}{0pt}\setlength{\parskip}{\baselineskip}\setstretch{1.0}\raggedright}
  {\par}

\makeatletter
\renewcommand{\maketitle}{%
  \begin{center}
  {\bfseries\@title\par}
  \vspace{1.5\baselineskip}
  {\@author\par}
  \end{center}
  \vspace{1\baselineskip}
}
\makeatother
